\section{Interoperability use case}
\new{HCM study metadata travel through exchange formats such as ISA \cite{isa}, and their statistics are described with OBI and STATO. HCMO does not import or align OBI and STATO as HCMO classes; the release shows at the instance level how HCMO data connect to that ecosystem, in two steps (Fig.~\ref{fig:evidence-layers}).}

\textbf{A recording-to-result evidence slice.} A pinned example, \texttt{examples/isa-hcmo-bridge.ttl}\fnurl{https://github.com/dhuzard/HCMO/tree/v0.3.0/examples}, uses PROV-O, specific OBI and STATO types, and ISA/Bioschemas exchange terms together with HCMO in one acyclic workflow from a raw recording to a statistical result \cite{isa,rocrate}. This is validated interoperability evidence, not an HCMO class mapping or a claim of formal ISA RO-Crate conformance.

\begin{figure}[t]
\centering
\begin{tikzpicture}[
  node distance=4mm and 6mm,
  box/.style={draw, rounded corners=1.5pt, align=center, font=\scriptsize, inner sep=3.5pt},
  flow/.style={-{Latex[length=1.8mm]}, font=\tiny}
]
\node[box, fill=orange!12] (sensor) {\texttt{hcm-tech:Sensor}\\DVC capacitive board\\12 electrodes, 4 Hz};
\node[box, fill=blue!8, right=of sensor] (obs) {\texttt{sosa:Observation}\\bin 11:00--11:01\\(real cohort 7623)};
\node[box, fill=green!9, right=of obs] (result) {\texttt{hcm-obs:QuantityValue}\\7.6 \texttt{unit:PERCENT}};
\node[box, fill=red!8, below=of obs] (group) {\texttt{hcm-bio:ExperimentalGroup}\\animals of cage 3401};
\node[box, fill=yellow!12, below=of sensor] (cage) {\texttt{hcm:MonitoredEnclosure}\\cage B6\_M\_12\_3401\\rack AAAA, position A6};
\node[box, fill=violet!8, below=of result] (series) {\texttt{hcm-tech:TimeSeries}\\CSV export\\attributed to DVC Analytics};
\draw[flow] (sensor) -- node[above]{madeBySensor} (obs);
\draw[flow] (obs) -- node[above]{hasResult} (result);
\draw[flow] (obs) -- node[right]{hasFeatureOfInterest} (group);
\draw[flow] (obs) -- node[left, pos=0.35]{occursIn} (cage);
\draw[flow] (sensor) -- node[left]{installedIn} (cage);
\draw[flow] (group) -- node[below]{housing assignment} (cage);
\draw[flow, dashed] (result) -- node[right]{row in} (series);
\end{tikzpicture}
\caption{Sensor $\rightarrow$ observation $\rightarrow$ result chain for one real
DVC activation-index bin. The cage-level signal has no individual subject, so the
feature of interest is the cage group; the derivation is proprietary, so only the
export file carries an attribution.}
\label{fig:dvcchain}
\end{figure}


\textbf{A 2 $\times$ 2 factorial study with repeated measures.} \new{The second step is a complete, generated study. Eight individually housed animals are allocated to a 2 $\times$ 2 factorial design --- treatment (vehicle or active) crossed with cage enrichment (standard or enriched) --- giving four groups of two. Each animal's dark-phase activity is observed on seven consecutive days, so day is a repeated, within-animal measure rather than a third crossed factor (56 observations in total), and one animal is re-housed during the study. A 2 $\times$ 2 design is the smallest one that exercises two crossed factors, their groups, and an interaction; the repeated measures and the re-housing test that observations stay attached to the right animal and enclosure over time. The fixture asks whether this structure --- animal identity, housing history, factors, groups, repeated observations, a derived tissue sample, and the statistical results --- survives conversion between HCMO and ISA, and what is lost when it does not.}

The fixture contains non-overlapping re-housing, an actual derived tissue \nolinkurl{bioschemas:Sample}, and separate fitted-model, estimate, confidence-interval, p-value, \nolinkurl{schema:MediaObject}, and file-fragment entities. Dark-phase outcomes use phenomenon intervals, and the re-housed animal's observations resolve to its new enclosure through the authoritative assignment. The declared mixed model is executed against the generated CSV with pinned numerical dependencies; the model serialization and the active-versus-vehicle contrast at standard enrichment use distinct exact row fragments. The input matrix (\nolinkurl{dark-phase-activity.csv}), the STATO-annotated result matrix (\nolinkurl{model-results.csv}), the STATO-typed entities and exact fragment links (\nolinkurl{canonical.ttl}), and an artifact guide (\nolinkurl{README.md}) are in \nolinkurl{examples/isa-roundtrip/} in the same folder.

\textbf{Three projections, two levels of fidelity.} Canonical HCMO RDF and an extended ISA RO-Crate JSON-LD serialisation are graph-isomorphic at 1,539 triples. \nolinkurl{roc-validator} 0.11.3 passes all RO-Crate 1.2 required rules and all ISA-specific required rules; its full inherited ISA run has one isolated upstream conflict, since the embedded ISA profile requires RO-Crate 1.1 although its minimal ISA fixture uses the 1.2 context. Formal ISA conformance therefore remains deferred: the draft profile has no permanent URI and its validator disagrees about the base RO-Crate edition. Native ISA-JSON and ISA-Tab are controlled-loss projections. Pinned \nolinkurl{isatools} 0.14.3 reproducibly generates checked-in ISA-JSON and ISA-Tab files and validates their JSON--Tab--JSON round trip for the genuine Source-to-Sample collection: a distinct animal Source, a real tissue Sample, their collection process, and their HCMO identifiers, which survive as ISA comments. Housing, direct whole-animal recording, source-bound factor assignments, explicit groups, repeated observations, and semantic STATO/file-fragment links remain declared losses, checked against loss manifests, rather than being represented with fabricated Sample proxies (Fig.~\ref{fig:roundtrip}). The generated files and modelling decisions are documented under \nolinkurl{native-isa/} in the same folder.

\begin{figure}[t]
\centering
\begin{tikzpicture}[
  x=0.93cm,
  y=0.93cm,
  box/.style={draw, rounded corners=1.5pt, align=center, font=\scriptsize,
    inner sep=3.5pt, text width=2.15cm},
  native/.style={box, draw=green!45!black, very thick, fill=green!8},
  extended/.style={box, draw=blue!55!black, fill=blue!7},
  context/.style={box, draw=orange!70!black, fill=orange!9},
  design/.style={box, draw=violet!60!black, fill=violet!7, text width=2.55cm},
  flow/.style={-{Latex[length=1.8mm]}, semithick},
  nativeflow/.style={flow, draw=green!45!black, very thick},
  contextflow/.style={flow, dashed, draw=orange!70!black},
  note/.style={font=\tiny, align=left}
]

% Experimental design and contextual records.
\node[design] (design) at (1.55,2.05)
  {$2\!\times\!2$ design\\treatment: vehicle/active\\enrichment: standard/enriched\\four groups; two animals each};
\node[native] (animal) at (5.35,2.05)
  {animal-8\\\textbf{ISA Source}\\HCMO Subject};
\node[context] (housing) at (9.25,2.05)
  {cage + time-bounded\\HCMO housing\\assignment; 9 rows};
\draw[flow] (design) -- (animal);
\draw[contextflow] (housing) -- (animal);

% Native Source-to-Sample branch: kept left of the acquisition branch so that
% no arrow crosses another node or edge.
\node[native] (sample) at (0.95,0)
  {tissue specimen\\\textbf{ISA Sample}};
\node[native] (collection) at (3.25,0)
  {collection\\Protocol + Process};
\draw[nativeflow] (animal.south west) -- (collection.north east);
\draw[nativeflow] (collection) -- (sample);

% Full HCMO / extended ISA RO-Crate acquisition and analysis branch.
\node[extended] (recording) at (7.65,0)
  {recording\\Protocol + Process\\56 observations};
\node[extended] (raw) at (10.35,0)
  {raw ISA DataFile\\\texttt{dark-phase-}\\\texttt{activity.csv}};
\draw[flow] (animal.south east) -- (recording.north west);
\draw[flow] (recording) -- (raw);

\node[extended] (results) at (7.65,-2.05)
  {STATO-annotated\\result matrix +\\exact row fragments};
\node[extended] (model) at (10.35,-2.05)
  {linear mixed-model\\Protocol + Process};
\draw[flow] (raw) -- node[right, font=\tiny]{input} (model);
\draw[flow] (model) -- (results);

\node[note, text width=11.3cm, anchor=west] at (0,-3.45) {%
  \textcolor{green!45!black}{\rule{4mm}{0.7pt}} also generated by native ISA-API;
  \textcolor{blue!55!black}{\rule{4mm}{0.7pt}} full RDF/extended-crate path;
  \textcolor{orange!70!black}{\rule[0.5ex]{4mm}{0.5pt}} contextual link, not data production.\\[-1pt]
  \texttt{examples/isa-roundtrip/}: \texttt{README.md}; \texttt{canonical.ttl};
  \texttt{ro-crate-metadata.json};\\[-1pt]
  \texttt{data/dark-phase-activity.csv}; \texttt{data/model-results.csv};
  \texttt{native-isa/}; \texttt{loss/}.};

\end{tikzpicture}
\caption{\new{Executable 2 $\times$ 2 factorial fixture with repeated measures (seven daily observations per animal).} Each entity and output has a distinct role. The green
animal Source $\rightarrow$ collection Process $\rightarrow$ genuine tissue
Sample path is reproducibly generated as ISA-JSON and ISA-Tab. The remaining
factor, housing, whole-animal recording, repeated-observation, and semantic
STATO/file-fragment structures are validated in HCMO RDF and the extended
ISA RO-Crate, but remain declared losses in the native ISA projection. The
result CSV exposes the reviewed STATO IRIs and labels while the RDF retains the
typed entities and exact row-fragment links.}
\label{fig:roundtrip}
\end{figure}

